\documentclass[a4paper,twoside,11pt]{book}
\usepackage[utf8]{inputenc}
\usepackage[T1]{fontenc}
\usepackage[english,german]{babel}

%%%%%%%%%%%%%%%%%%%%%%%%%%%% NAME AND TITLE %%%%%%%%%%%%%%%%%%%%%%%%%%%%%%%%%%%%%%%%%%%%
%% Bitte füllen Sie dieses mit ihren Daten aus.
\title{Das Universum und der ganze Rest}
\author{B.Sc.Max Mustermann}
\newcommand{\firstprof}{Prof. Dr.-Ing. G. von Voigt}
\newcommand{\secondprof}{Prof. Dr. rer. nat. M. Smith}
\newcommand{\betreuer}{Dipl.-Inform. H. Tobaben}
\date{\today}
\newcommand{\art}{Masterarbeit}  % Bachelorarbeit / Diplomarbeit / Studienarbeit
\newcommand{\studiengang}{Informatik (M.Sc.)}
\newcommand{\keywords}{Security,Masterarbeit,Wichtig,Test}
%%%%%%%%%%%%%%%%%%%%%%%%%%%%%%%%%%%%%%%%%%%%%%%%%%%%%%%%%%%%%%%%%%%%%%%%%%%%%%%%%%%%%%%%

% Trennungsregeln
\usepackage{hyphenat}
\hyphenation{Grid--Informations-systeme Sche-du-ler Sche-du-ling Per-for-mance Meta-sche-duler  Pla-gi-ats-prü-fung}

\usepackage{amssymb}							% provides mathmatical math symbols such as dots, arrows, etc.
\usepackage{amsmath}							% provides align environment
\usepackage{listings}
\lstset{captionpos=b,frame=single}
\usepackage[german]{fancyref}
%% for fancy refereing to listings.
\newcommand*{\fancyreflstlabelprefix}{lst}

\fancyrefaddcaptions{english}{%
  \providecommand*{\freflstname}{listing}%
  \providecommand*{\Freflstname}{Listing}%
}

\fancyrefaddcaptions{german}{%
  \providecommand*{\freflstname}{listing}%
  \providecommand*{\Freflstname}{Listing}%
}

\frefformat{plain}{\fancyreflstlabelprefix}{\freflstname\fancyrefdefaultspacing#1}
\Frefformat{plain}{\fancyreflstlabelprefix}{\Freflstname\fancyrefdefaultspacing#1}

\frefformat{vario}{\fancyreflstlabelprefix}{%
  \freflstname\fancyrefdefaultspacing#1#3%
}
\Frefformat{vario}{\fancyreflstlabelprefix}{%
  \Freflstname\fancyrefdefaultspacing#1#3%
}

%% 1.5 facher Zeilenabstand.
\usepackage[onehalfspacing]{setspace}

\makeatletter
\let\thetitle\@title
\let\theauthor\@author
\let\theday\@date
\makeatother

\usepackage{geometry}
\geometry{a4paper,left=35mm,right=25mm, top=2.5cm, bottom=2.5cm}

% Metadaten
\usepackage[
	colorlinks=true,
	linkcolor=black,						% enable for printing!
	urlcolor=black,						  % enable for printing!
	citecolor=black,						% enable for printing!
	%pdfstartview=Fit,						% fits the page to the window; other options: FitH/FitV (horiz./vert.)...
	%pdfpagelayout=TwoPageLeft,		% the way pages are displayed, options: SinglePage, OneColumn, TwoColumnLeft|Right, TwoPageLeft|Right
	final=true,									% turn on all processing options
	plainpages=false,						% Forces page anchors to be named by the arabic form of the page number, rather than the formatted form.
	pdfpagelabels]							% set PDF page labels
	{hyperref}
\hypersetup{
  pdftitle={\thetitle},
  pdfauthor={\theauthor},
  pdfsubject={\keywords}
}

\usepackage[nopostdot,nonumberlist,acronym]{glossaries}
%\makeglossaries

\makenoidxglossaries
%%%%%%%%%%%%%%%%%%%%%%%%%%%%%%%%%%%%%%%%%%%%%%%%%%%%%%%%%%%%%
%%                                                         %%
%%            Fügen sie hier alle Abkürzungen ein,         %%
%%                 welche Sie benutzen                     %%
%%                                                         %%
%%%%%%%%%%%%%%%%%%%%%%%%%%%%%%%%%%%%%%%%%%%%%%%%%%%%%%%%%%%%%

%% Name für Befehl \gls{} / Abkürzung / Ausgeschrieben
\newacronym{dcsec}{DCSec}{Distributed Computing \& Security Research Group}
\newacronym{mpi}{MPI}{Message Passing Interface}
\newacronym{omp}{OpenMP}{Open Multi-Processing}
\newacronym{pag}{PAG}{Per Anhalter durch die Galaxis}
\newacronym{xml}{XML}{eXtendable Markup Language}

%% Mehr informationen können unter https://www.ctan.org/pkg/glossaries?lang=de gefunden werden.

\usepackage{graphicx}
% sorgt dafür, dass LaTeX in images nach den Bildern sucht und der Ordner "etwas" aufgeräumter ist.
\graphicspath{ {./images/} {./graphics/} }


\usepackage{fancyhdr}
\pagestyle{fancy}
\fancyhf{} % sets both header and footer to nothing
\renewcommand{\headrulewidth}{0pt}
\parindent 0pt

\usepackage{helvet}
\renewcommand{\familydefault}{\sfdefault}
\fontfamily{phv}\selectfont


%%%%%%%%%%%%%%%%%%%%%%%%%%%%%%%%%%%%%%%%%%%%%%%%%%%%%%%%%%%%%%%%%%%%%%%%%%%%%%%%%%%%%%%%
%% Falls Sie noch weitere Pakete benötigen können diese hier eingebunden werden.

%\usepackage{...}

%%%%%%%%%%%%%%%%%%%%%%%%%%%%%%%%%%%%%%%%%%%%%%%%%%%%%%%%%%%%%%%%%%%%%%%%%%%%%%%%%%%%%%%%

\setlength{\headheight}{13.6pt}
\begin{document}

\frontmatter
%%%%%%%%%%%%%%%%%%%%%%%%%%%%%%%%%%%%%%%%%%%%%%%%%%%%%%%%%%%%%
%%                                                         %%
%%    Hier sollten von ihrer Seite aus keine Änderungen    %%
%%                  notwendig sein.                        %%
%%                                                         %%
%%%%%%%%%%%%%%%%%%%%%%%%%%%%%%%%%%%%%%%%%%%%%%%%%%%%%%%%%%%%%

\newgeometry{left=3cm,right=4cm, top=3cm, bottom=3.5cm}
\begin{titlepage}

%
% Logos
%
\includegraphics[height=1cm]{logo_dcsec} \hfill
\includegraphics[height=1cm]{logo_luh}

\vspace{-3mm}

\rule{\textwidth}{1pt}

\vspace{5mm}

\large

Gottfried Wilhelm Leibniz Universität Hannover

Institut für Verteilte Systeme

Distributed Computing \& Security Group

\vspace{4.0cm}

\art\\
im Studiengang \studiengang

\vspace{1.0cm}
\huge{\thetitle}

\vfill


\large
\begin{tabular}{p{3cm}l}
Verfasser:    & \theauthor\\
Betreuer:     & \betreuer\\
Datum:        & \theday
\end{tabular}


\end{titlepage}

\fancyfoot[C]{\thepage}

\tableofcontents

%% sorgt dafür, dass Verzeichnisse nicht auf ungeraden Seiten erscheinen, sondern nacheinander.
\makeatletter
\@openrightfalse
\makeatother

% Hier kommt das Abbildungsverzeichnis.
\listoffigures

% Hier kommt eine Liste der Tabellen.
\listoftables

% Codeverzeichnis
\renewcommand*{\lstlistlistingname}{Codeverzeichniss}
\lstlistoflistings

\printnoidxglossaries

%\printglossary[title=Nomenclature]

\mainmatter
\renewcommand{\headrulewidth}{0.4pt}
\fancyhead[OR]{\fontfamily{phv}\selectfont \rightmark}
\fancyhead[EL]{\fontfamily{phv}\selectfont \leftmark}

\makeatletter
\@openrighttrue
\makeatother

%%%%%%%%%%%%%%%%%%%%%%%%%%%%%% HAUPTTEIL %%%%%%%%%%%%%%%%%%%%%%%%%%%%%%%%%
%% binden sie hier Ihre anderen LaTeX Dateien ein.
\chapter{Einleitung}
\label{chap:Einleitung}

\section{Handbuch}
\label{sec:Handbuch}
%% Das Label sollte etwas über den Typ aussagen. Section-Labels sollten mit sec, Chapters mit chap, Bilder mit fig und Tabellen mit tab beginnen. So kann fref den Typ erkennen und entsprechend referenzieren.
Der Reiseführer \gls{pag} hat eine sehr schöne Theorie zur Komplexität des Universums:

\begin{quote}
Es gibt eine Theorie, die besagt, wenn jemals irgendwer genau herausfindet, wozu das Universum da ist und warum es da ist, dann verschwindet es auf der Stelle und wird durch noch etwas Bizarreres und Unbegreiflicheres ersetzt. - Es gibt eine andere Theorie, nach der das schon passiert ist. \hfill \cite{adams1995hitchhiker}
\end{quote}

Sie ist natürlich nicht so schön wie das Logo unserer Uni in \Fref{fig:uni-logo} oder das Logo der \gls{dcsec} in \Fref{fig:dcsec-logo}.

\begin{figure}[ht]
\centering
\includegraphics[width=\textwidth]{logo_luh}
\caption[Uni-Logo]{Das Uni-Logo}
\label{fig:uni-logo}
\end{figure}

\begin{figure}[ht]
\centering
\includegraphics{logo_dcsec}
\caption[DCSec-Logo]{Distributed Computing Security Group Logo}
\label{fig:dcsec-logo}
\end{figure}

\begin{table}
\begin{center}
\begin{tabular}{lcr}
  Spalte 1 & Spalte 2 & Spalte 3 \\
  10 & 20 & 30 \\
\end{tabular}
\end{center}
\label{tab:test}
\caption{Testtabelle}
\end{table}

Was man vielleicht noch beachten sollte ist, dass sind Abkürzungen wie \gls{dcsec} oder \gls{pag} bereits einmal benutzt worden. Werden sie zukünftig nicht mehr neu ausgeschrieben, im Gegensatz zu unbenutzten wie \gls{xml}
\chapter{Weltherrschaft}
\label{chap:weltherrschaft}


Der Weg zur Weltherrschaft ist ganz einfach wie \Fref{lst:weltherrschaft} zeigt.

\begin{lstlisting}[label=lst:weltherrschaft, caption={How to reach World domination}]
while not reached
    try reach world domination
end
\end{lstlisting}
\input{content/8_ende}


%%%%%%%%%%%%%%%%%%%%%%%%%%%%%%%%%%%%%%%%%%%%%%%%%%%%%%%%%%%%%%%%%%%%%%%%%%


\appendix


%%%%%%%%%%%%%%%%%%%%%%%%%%%%%%% ANHÄNGE %%%%%%%%%%%%%%%%%%%%%%%%%%%%%%%%%%
%% binden sie hier Ihre LaTeX Dateien ein, welche zum Anhang gehören.

%%%%%%%%%%%%%%%%%%%%%%%%%%%%%%%%%%%%%%%%%%%%%%%%%%%%%%%%%%%%%%%%%%%%%%%%%%

\bibliographystyle{plain}
\bibliography{referenzen}

\end{document}
